\section{Discussion}
\label{sec:discussion}

\subsection{From Single-Day to Multi-Day: Bridging Two Temporal Scales}

The progression from single-day detection (71.5\%) to multi-day regime identification (81.2\% in 2024) reveals complementary validation dimensions. Single-day testing establishes that the LLM understands dealer mechanics within individual trading sessions; regime detection validates that this understanding extends to persistent multi-day structural states. The raw chain finding (92.3\% detection from first-principles data) provides additional support: if the model merely matched statistical signatures, raw chain performance would degrade rather than improve when pre-calculated summaries are removed.

The detection-alpha orthogonality finding from single-day analysis (Sharpe collapse from 1.8 to 0.1 while detection remains stable) motivated the regime extension: patterns that persist despite economic insignificance must reflect structural mechanics rather than exploitable inefficiencies. This distinction---mechanism identification versus alpha generation---is fundamental to deploying LLMs for market surveillance and risk management rather than trading.

\subsection{Validating Structural Reasoning}

Five-phase validation provides evidence that the LLM performs structural reasoning rather than pattern matching. The 69.1 percentage point discrimination between 2024 (81.2\%) and 2020 (12.1\%) establishes selectivity: if the model relied on memorized associations or universal statistical patterns, detection rates would converge across periods rather than exhibiting 5.7$\times$ separation.

Three findings support this interpretation. First, the LLM maintained 98\% mechanical accuracy across both detection conditions---correctly computing persistence, magnitude, and flip counts regardless of whether windows qualified as regimes. Second, negative controls achieved 0\% false positives on transitional and low-magnitude synthetic data, confirming the model enforces structural criteria rather than over-detecting patterns. Third, confidence scores provide continuous quality discrimination beyond binary classification ($r > 0.4$ with regime metrics), suggesting the model tracks regime robustness, not just threshold crossing.

This selectivity contrasts with earlier 5-day trajectory testing \citep{regan2025obfuscation}, which achieved 98--100\% detection across all market conditions by identifying universal daily hedging flows. The deliberate shift to 30-day windows with strict criteria transformed the task from universal pattern matching to selective regime identification---a critical distinction for validating genuine structural reasoning.

\subsection{Market Structure Evolution and 0DTE Hypothesis}

The multi-year temporal analysis (Phase 5) reveals gradual regime evolution tracking 0DTE options adoption rather than sharp structural breaks. Detection progressed from 12.2\% (2020) through borderline years (3.7\% in 2021, 32.4\% in 2022) to sustained 100\% detection in 2024--2025, with GEX magnitude growing 360\% (\$5B $\rightarrow$ \$23B).

The non-monotonic pattern (32.4\% in 2022 dipping to 20.2\% in 2023 before reaching 100\% in 2024) is informative: 2023's elevated volatility (FOMC uncertainty, banking stress) repeatedly disrupted regime formation despite growing 0DTE hedging pressure, suggesting regime persistence requires both sustained dealer pressure \textit{and} a volatility environment permitting consolidation. This tipping-point dynamic strengthens the structural interpretation.

While the temporal coincidence with 0DTE proliferation is consistent with a causal relationship, we acknowledge that concurrent factors---interest rate changes (0.25\% $\rightarrow$ 5.5\%), volatility regime shifts, quantitative trading growth, and market maker concentration---may contribute. However, the non-monotonic detection pattern argues against gradual secular trends as primary drivers: detection remained flat through 2023 despite continuous interest rate increases and technology diffusion, then jumped discontinuously in 2024 coinciding with 0DTE market saturation ($\approx$48\% volume share).

\subsection{Temporal Obfuscation as Generalizable Methodology}

Our obfuscation testing framework generalizes beyond financial applications. The core principle---stripping temporal and contextual identifiers to force structural reasoning---addresses a universal challenge in deploying LLMs for domain-specific analysis: distinguishing genuine reasoning from training data contamination. The methodology requires three components: (1) identifiable features that can be systematically removed, (2) structural patterns that persist after obfuscation, and (3) negative controls validating discrimination.

Potential applications include medical time series analysis (removing patient identifiers and dates while preserving physiological patterns), climate data analysis (obfuscating location and temporal markers), and industrial anomaly detection (stripping equipment identifiers). In each case, temporal obfuscation could validate whether LLMs reason from structural properties or rely on memorized associations from training data.

\subsection{Dispersed Knowledge and Information Aggregation}

The raw chain superiority finding (92.3\% vs.\ 61.5\%) admits a theoretical interpretation grounded in \citet{hayek1945}'s analysis of dispersed knowledge in economic systems. Hayek argued that the economic problem is fundamentally one of utilizing knowledge ``not given to anyone in its totality''---that relevant information is distributed across countless individual participants, each possessing local knowledge inaccessible through centralized aggregation. Traditional parametric metrics like net GEX attempt precisely such centralized aggregation: compressing the full strike-level options surface into a single scalar value. Our results suggest this compression is lossy in a theoretically predictable way.

The raw options chain preserves what Hayek termed ``knowledge of the particular circumstances of time and place''---the asymmetrical strike concentrations, implied volatility skew patterns, and spatial open interest distributions that reflect actual dealer positioning at each strike. When the LLM processes raw chain data, it functions as an information aggregator that synthesizes these dispersed local signals into a coherent structural assessment, identifying emergent order arising from countless decentralized dealer hedging constraints. The 30.8 percentage point improvement over GEX-assisted detection provides empirical evidence that this ``local knowledge'' embedded in individual strikes contains structural signal that scalar aggregation necessarily destroys.

This interpretation extends beyond methodological curiosity. The alpha disappearance puzzle---stable detection (68--74\%) despite collapsing profitability (Sharpe 1.8 $\rightarrow$ 0.1)---parallels a core insight from Austrian capital theory: that structural knowledge of market coordination patterns differs fundamentally from actionable entrepreneurial knowledge \citep{kirzner1973competition}. Detecting that dealers are constrained to pro-cyclical hedging is structural knowledge; profiting from that detection requires entrepreneurial judgment about timing, magnitude, and competing market forces that our framework deliberately excludes. The orthogonality between detection and alpha is thus not a limitation but a feature: it confirms the framework identifies genuine structural mechanics rather than ephemeral trading opportunities.

\subsection{Practitioner Implications}

For financial practitioners, our findings suggest four actionable insights. First, providing LLMs with raw structured data rather than pre-calculated summaries yields substantially better structural detection (92.3\% vs.\ 61.5\%), challenging common pipeline designs that aggregate data before analysis---an empirical vindication of preserving informational granularity over computational convenience. Second, the detection-alpha orthogonality result cautions against using structural pattern detection directly for trading; the value lies in risk management and market surveillance applications where mechanism identification matters more than profitability. Third, the 0DTE-driven regime shift suggests that market microstructure monitoring should evolve with product innovation---static models calibrated to pre-2022 data will miss the structural reorganization that our framework detects. Fourth, the dispersed knowledge interpretation suggests that financial AI systems may achieve better performance across domains by ingesting raw, granular data rather than pre-processed summaries---a design principle applicable to credit risk, fixed income analysis, and other areas where scalar metrics are standard.

\subsection{Limitations}

Six limitations merit discussion.

\noindent\textbf{Single asset}: Testing focused on SPY; cross-asset generalization (individual equities, other ETFs) remains untested. SPY's unique liquidity and 0DTE availability may limit transferability.

\noindent\textbf{End-of-day measurement}: Our GEX\_OI approach captures dealer inventory but not intraday gamma dynamics; high-frequency flow data could refine detection, particularly for 0DTE contracts that expire within a single trading session.

\noindent\textbf{Causal attribution}: The 0DTE hypothesis is supported by temporal coincidence and theoretical mechanism but remains circumstantial; observational data cannot exclude alternative explanations. A natural experiment (e.g., 0DTE suspension) would provide stronger causal evidence.

\noindent\textbf{Single LLM}: All experiments use OpenAI o4-mini; cross-model validation with alternative reasoning models would strengthen generalizability claims.

\noindent\textbf{Shuffle test anomaly}: The 61\% false positive rate on 2024 shuffled data (against a $<$20\% criterion) reflects extreme regime persistence rather than framework failure---2024 regimes are so dominant that random permutation rarely disrupts the aggregate signal. This is documented as a known finding rather than treated as a validation failure.

\noindent\textbf{Threshold sensitivity}: All 25 tested parameter combinations maintained $>$72pp discrimination, but thresholds represent empirically validated design choices rather than first-principles derivations. Future work should explore adaptive thresholding based on market conditions.
